\documentclass[12pt,a4paper]{article}

% ── Packages ──────────────────────────────────────────────
\usepackage[margin=1in]{geometry}
\usepackage{times}
\usepackage{setspace}
\usepackage{graphicx}
\usepackage{enumitem}
\usepackage{hyperref}
\usepackage{fancyhdr}
\usepackage{array}

\hypersetup{colorlinks=false}

% ── Spacing ───────────────────────────────────────────────
\onehalfspacing
\setlength{\parindent}{0pt}
\setlength{\parskip}{8pt}

\pagestyle{fancy}
\fancyhf{}
\rfoot{\thepage}
\renewcommand{\headrulewidth}{0pt}

% ══════════════════════════════════════════════════════════
\begin{document}

% ── Title Page ────────────────────────────────────────────
\begin{titlepage}
    \centering
    \vspace*{0.5cm}

    {\large\bfseries CHRIST (Deemed to be University), Bengaluru}\\[4pt]
    {\normalsize Department of Computer Science}\\[1.5cm]

    {\LARGE\bfseries Project Synopsis}\\[0.4cm]
    {\large Course: CSC481-5 --- Computer Science Project}\\[0.3cm]
    {\large Semester VI}\\[2cm]

    {\huge\bfseries ScholarLens}\\[0.3cm]
    {\Large Offline SLM-Based Research Literature Review Agent}\\[2.5cm]

    \begin{tabular}{l l l}
        \textbf{Sl. No.} & \textbf{Student Name} & \textbf{Register No.} \\[6pt]
        1. & Aditya Mehta   & 2440204 \\[4pt]
        2. & Jitvan Chadha  & 2440222 \\[4pt]
        3. & Ritesh K R     & 2440233 \\[4pt]
    \end{tabular}

    \vspace{2cm}

    \begin{tabular}{>{\bfseries}r l}
        Guide: & Dr. New Begin \\[4pt]
        Class: & 5BSC CS \\[4pt]
    \end{tabular}

    \vfill
    {\large June 2026}
\end{titlepage}

% ── Synopsis (Abstract) ──────────────────────────────────

\begin{center}
    {\Large\bfseries Synopsis}
\end{center}

\vspace{0.3cm}

\textbf{Title:} ScholarLens --- An Offline SLM-Based Research Literature
Review Agent

\textbf{Problem Statement:}
Conducting a thorough literature review is a critical yet time-intensive
task for researchers and students.  The process involves reading,
organising, and synthesising large numbers of published papers to
understand the state of the art in a given field.  While several
AI-powered tools (e.g.\ Elicit, Semantic Scholar) have emerged to assist
with this process, they universally depend on cloud infrastructure and
internet connectivity.  This raises concerns around data privacy,
accessibility in low-connectivity environments, and recurring costs that
are prohibitive for students.  There is currently no integrated, offline
tool that combines document ingestion, intelligent search, automated
summarisation, and research gap identification in a single
privacy-preserving application.

\textbf{Proposed Solution:}
ScholarLens is a fully offline, privacy-first research literature review
agent powered by locally-running Small Language Models (SLMs).  The system
allows users to upload research papers in PDF format, which are then
parsed, chunked, and embedded into a local vector database.  Users can
search their paper collection using a hybrid approach that combines
semantic similarity with keyword-based retrieval.  The system further
provides SLM-driven capabilities including paper summarisation, key claim
extraction, thematic clustering, and research gap analysis---all executed
entirely on the user's machine without any data leaving the local
environment.

\textbf{Technology Stack:}
The backend is developed in Python using the FastAPI framework.
The frontend is a React.js single-page application built with Vite.
Language model inference is handled by Ollama running locally, while
ChromaDB serves as the persistent vector database.
Sentence-Transformers provide the embedding generation, and BM25 is used
for keyword-based retrieval.  PDF processing utilises PyMuPDF and
pdfplumber.

\textbf{Modules:}
\begin{enumerate}[leftmargin=2em, nosep]
    \item \textbf{Document Ingestion} --- PDF parsing, text chunking, and
          metadata extraction.
    \item \textbf{Knowledge Base} --- Embedding generation and vector
          storage via ChromaDB.
    \item \textbf{Search \& Retrieval} --- Hybrid semantic + keyword search
          with ranked results.
    \item \textbf{Analysis} --- SLM-powered summarisation, claim extraction,
          and theme clustering.
    \item \textbf{Gap Finder} --- Research gap identification and direction
          suggestion.
\end{enumerate}

\textbf{Expected Outcome:}
A functional, self-contained web application that enables researchers and
students to manage, search, analyse, and derive insights from their
research paper collection entirely offline, ensuring complete data privacy
and zero recurring costs.

\end{document}
